\documentclass[11pt, letterpaper]{article}

\usepackage[margin=0.9in]{geometry}
\usepackage{booktabs}
\usepackage{graphicx}
\usepackage{hyperref}
\usepackage{amsmath, amssymb}
\usepackage{enumitem}
\usepackage{xcolor}
\usepackage{titlesec}
\usepackage{fancyhdr}
\usepackage{multirow}
\usepackage{caption}
\usepackage{tcolorbox}
\usepackage{float}
\usepackage{array}
\usepackage{pifont}

% NYT Connections colors
\definecolor{yellow_conn}{HTML}{F9DF6D}
\definecolor{green_conn}{HTML}{A0C35A}
\definecolor{blue_conn}{HTML}{B0C4EF}
\definecolor{purple_conn}{HTML}{BA81C5}
\definecolor{codebg}{HTML}{F5F5F0}
\definecolor{highlightbg}{HTML}{FFF8DC}
\definecolor{accent}{HTML}{1A5490}

\hypersetup{colorlinks=true, linkcolor=accent, urlcolor=accent}

\pagestyle{fancy}
\fancyhf{}
\rhead{STA 561D | Duke}
\lhead{Infinite Connections --- Executive Summary}
\rfoot{\thepage}
\renewcommand{\headrulewidth}{0.4pt}

\titleformat{\section}{\Large\bfseries\color{accent}}{\thesection.}{0.5em}{}
\titleformat{\subsection}{\large\bfseries}{\thesubsection}{0.6em}{}

\newtcolorbox{keybox}{
    colback=highlightbg,
    colframe=yellow!50!orange,
    arc=2mm,
    boxrule=0.5pt,
    left=6pt,
    right=6pt,
    top=4pt,
    bottom=4pt,
}

\title{\textbf{Infinite Connections}\\[0.3em]
\large An AI System for Generating NYT-Style Connections Puzzles\\[0.3em]
\normalsize Executive Summary of Current Results and Implementation Pipelines (v2)}
\author{STA 561D --- Duke University}
\date{April 2026}

\begin{document}

\maketitle
\thispagestyle{fancy}

\vspace{-1em}

%-----------------------------------------------------------------------------
\section{Project Overview}

The New York Times Connections puzzle presents players with 16 words that must be
partitioned into 4 groups of 4 based on a hidden theme. \textbf{Infinite Connections}
is a system that generates such puzzles \emph{automatically} at scale, validates them
with multiple independent solvers, and serves them through a playable web interface.

We have built and benchmarked \textbf{two complete pipelines}:

\begin{enumerate}[leftmargin=*, itemsep=0.2em]
\item \textbf{Pipeline A --- LLM-Heavy (baseline)}: 5 LLM calls per puzzle; LLM
proposes word pools, MPNET embeddings refine selection.
\item \textbf{Pipeline B --- Category-First Retrieval (CFR, novel)}, two modes:
    \begin{itemize}[itemsep=0.1em]
    \item \textbf{Mode A (remix, 0 LLM calls)}: samples categories from the
    2{,}054-category NYT pool, retrieves words via KNN.
    \item \textbf{Mode B (fresh, 1 LLM call)}: LLM writes 4 fresh categories in one
    batched call; retrieves words via KNN.
    \end{itemize}
\end{enumerate}

\textbf{CFR v2 (rubric-compliant upgrade).} A careful re-read of the rubric surfaced
two constraints the first-pass CFR did not meet. We addressed both:
\begin{itemize}[leftmargin=*, itemsep=0.1em]
\item \emph{``If you generate a past connections puzzle you will automatically fail.''}
We added an exact 4-word-group match check against all ${\sim}2{,}216$ past NYT groups;
if the top-scoring retrieval would reproduce an NYT group, the next-best 4-word
combination is used instead.
\item \emph{Professor guidance: the word bank should be larger than the original
game's.} We augmented the 4{,}918-word NYT bank with 10{,}000+ common English lemmas
from WordNet, yielding a \textbf{14{,}877-word bank (3$\times$ larger than NYT's)}.
\end{itemize}

In addition, we deployed a \textbf{playable web application} hosted via GitHub Pages
that loads 542 validated puzzles and replicates the NYT Connections UI.

\begin{keybox}
\textbf{Headline result (v2).} Pipeline B (CFR) achieves a \textbf{99\% validator pass
rate} while drawing \textbf{62--69\% of its output words from non-NYT sources},
reproducing \textbf{zero past NYT groups} across 200 benchmark puzzles, and running at
\textbf{1/20 the cost} and \textbf{4$\times$ the speed} of the baseline LLM-heavy
pipeline.
\end{keybox}

%-----------------------------------------------------------------------------
\section{Pipeline A --- LLM-Heavy Baseline}

Pipeline A follows the architecture from ``Making New Connections'' (arXiv:2407.11240).
For each puzzle it issues 5 LLM calls:

\begin{enumerate}[leftmargin=*, itemsep=0.15em]
\item \textbf{Calls 1--4 (Group creation).} Four iterative calls, one per group. Each
returns a category name plus 8 candidate words. Story injection (4 random seed words
from the NYT bank) is embedded in the system prompt to force diversity.
\item \textbf{MPNET selection (non-LLM).} Of the 8 candidates, the best 4 are chosen
by enumerating all $\binom{8}{4}=70$ subsets and picking the one with the highest
average pairwise cosine similarity.
\item \textbf{Call 5 (Editor pass).} A final LLM call reviews all 4 category names
and rewrites any that don't accurately describe their words.
\item \textbf{Difficulty color assignment (non-LLM).} The 4 groups are sorted by
their avg pairwise cosine similarity; highest $\rightarrow$ yellow, lowest $\rightarrow$ purple.
\item \textbf{Deduplication (non-LLM).} The 16-word set is checked against all
554 existing NYT puzzles; flagged if $>\!6$ words overlap with any single puzzle.
\item \textbf{Roundtable validation (non-LLM).} Two independent solvers
(an embedding solver and a beam-search clustering solver) attempt to recover the
intended grouping. A puzzle is accepted if either solver fully solves it.
\end{enumerate}

\subsection*{Pipeline A Results}

\begin{table}[H]
\centering
\small
\begin{tabular}{lrrrr}
\toprule
\textbf{Batch} & \textbf{Model} & \textbf{Count} & \textbf{Valid} & \textbf{Pass rate} \\
\midrule
100-puzzle (gpt-4o-mini) & gpt-4o-mini & 100 & 91 & 91\% \\
100-puzzle (gpt-4o)      & gpt-4o      & 100 & 87 & 87\% \\
575-puzzle batch         & gpt-4o-mini & 575 & 534 & 93\% \\
\midrule
\textbf{Total}           &             & \textbf{775} & \textbf{712} & \textbf{91.9\%} \\
\bottomrule
\end{tabular}
\caption{Pipeline A generation results. Total cost $\approx$\$2.00.}
\end{table}

\textbf{Observation:} gpt-4o produces more creative wordplay categories but lower
pass rate (87\%) because embedding-based validators struggle with pure wordplay.
gpt-4o-mini is a better fit for this pipeline.

%-----------------------------------------------------------------------------
\section{Pipeline B --- Category-First Retrieval (CFR)}

The key insight: in Pipeline A, the LLM wastes tokens \emph{generating candidate
words}, only to have MPNET pick 4 of them anyway. CFR flips the flow:

\begin{enumerate}[leftmargin=*, itemsep=0.15em]
\item \textbf{Propose 4 categories} (two modes):
    \begin{itemize}[itemsep=0.1em]
    \item \textbf{Mode A (``remix'')}: sample 4 mutually dissimilar categories from
    the 2,054 existing NYT categories (diversity enforced via pairwise MPNET cosine
    distance $\geq 0.35$). Zero LLM calls.
    \item \textbf{Mode B (``fresh'')}: \emph{one} batched LLM call returns 4 fresh
    category names in a single response.
    \end{itemize}
\item \textbf{Retrieve words per category.} For each category name, encode it with
MPNET and run KNN over the precomputed 4,918-word embedding matrix to get the top-30
semantically closest words. Filter morphological duplicates (TIME/TIMES/TIMER) and
category-token overlaps (``TIME'' for category ``TIME PERIODS'').
\item \textbf{Select best 4.} Reuse Pipeline A's MPNET selector to pick the 4 most
cohesive words from the top 8 candidates.
\item \textbf{Assign colors, deduplicate, validate.} Identical to Pipeline A.
\end{enumerate}

\subsection*{Offline precomputation (one-time, $\sim$5 minutes)}

\begin{itemize}[leftmargin=*, itemsep=0.1em]
\item Encode all 4{,}918 NYT words $\rightarrow$ $4{,}918 \times 768$ matrix.
\item Encode all 2{,}054 unique NYT categories $\rightarrow$ $2{,}054 \times 768$ matrix.
\item Build a scikit-learn \texttt{NearestNeighbors} index with cosine distance.
\item Cache the result as \texttt{data/cache/nyt\_embeddings.npz}.
\end{itemize}

\subsection*{Pipeline B Results (v1 vs. v2)}

\begin{table}[H]
\centering
\small
\begin{tabular}{lrrrrr}
\toprule
\textbf{Version \& Mode} & \textbf{Word bank} & \textbf{Pass rate} & \textbf{Non-NYT words} & \textbf{NYT-group collisions} & \textbf{Time/puzzle} \\
\midrule
v1 --- Mode A (remix) & 4{,}918 & 97\% & 0\% & not checked & 1.41 s \\
v1 --- Mode B (fresh) & 4{,}918 & 99\% & 0\% & not checked & 2.67 s \\
v2 --- Mode A (remix) & \textbf{14{,}877} & \textbf{99\%} & \textbf{62.6\%} & \textbf{0 / 100} & 1.27 s \\
v2 --- Mode B (fresh) & \textbf{14{,}877} & \textbf{99\%} & \textbf{69.1\%} & \textbf{0 / 100} & 2.77 s \\
\bottomrule
\end{tabular}
\caption{CFR pipeline benchmarks (100 puzzles per row). v2 expands the bank via WordNet
augmentation and enforces the rubric's ``no-past-puzzle'' rule via 4-word-group
matching.}
\end{table}

%-----------------------------------------------------------------------------
\section{Head-to-Head Comparison}

\begin{table}[H]
\centering
\small
\begin{tabular}{l r r r}
\toprule
\textbf{Metric} & \textbf{Pipeline A} & \textbf{CFR v2 Mode A} & \textbf{CFR v2 Mode B} \\
                & (baseline)          & (remix, 0 calls)        & (fresh, 1 call) \\
\midrule
LLM calls / puzzle              & 5                & \textbf{0}             & \textbf{1} \\
Cost / 1{,}000 puzzles          & \textasciitilde\$1.00 & \textbf{\$0.00}   & \textbf{\$0.05} \\
Wall time / puzzle              & \textasciitilde10 s    & \textbf{1.27 s}    & \textbf{2.77 s} \\
Pass rate                       & 87--93\%         & \textbf{99\%}          & \textbf{99\%} \\
Word bank size                  & 4{,}918 (NYT)    & \textbf{14{,}877}      & \textbf{14{,}877} \\
Non-NYT words in output         & 0\%              & \textbf{62.6\%}        & \textbf{69.1\%} \\
Hard-fail safety check          & word overlap only & \textbf{+ 4-word group}& \textbf{+ 4-word group} \\
NYT group collisions (of 100)   & not checked      & \textbf{0}              & \textbf{0} \\
API fragility                   & High             & \textbf{None}           & Low \\
\bottomrule
\end{tabular}
\caption{Pipeline comparison on 100-puzzle batches (CFR v2).}
\end{table}

\textbf{Speed:} CFR Mode A is $\sim$7$\times$ faster and Mode B is $\sim$4$\times$ faster
than the baseline, because the embedding retrieval (KNN lookup on 4,918 words) is
essentially instantaneous compared to network round-trips to an LLM API.

\textbf{Cost:} For a run of 10{,}000 puzzles, the baseline costs $\sim$\$10; CFR Mode A
costs \$0 and Mode B costs $\sim$\$0.50. A \textbf{20$\times$ cost reduction} at
\textbf{higher quality}.

\textbf{Methodological contribution:} The project's original contribution shifts from
``we called the LLM to generate puzzles'' to ``we built an embedding-first hybrid
pipeline where the LLM is a small, targeted component.''

%-----------------------------------------------------------------------------
\section{Rubric-Compliance Upgrade (CFR v1 $\rightarrow$ v2)}

The original CFR pipeline worked well by statistical metrics but missed two
rubric-level constraints. The v2 upgrade keeps the architecture identical and adds
two targeted safeguards.

\subsection*{Fix 1: Augmented word bank (new \texttt{src/cfr/word\_bank.py}).}

The professor noted that ``the total word bank should be more than that of the original
game's.'' We union three sources and deduplicate:

\begin{enumerate}[leftmargin=*, itemsep=0.1em]
\item The 4{,}918 unique NYT words (retained for NYT-style familiarity).
\item NLTK WordNet single-token lemmas with nonzero tagged-corpus frequency --- filtered
to 3--15 chars, alphabetic only, no profanity. This yields $\sim$10{,}000 common
English words (no archaic/technical noise).
\item An optional Google 10k most-common-English list if present (not required).
\end{enumerate}

Result: \textbf{14{,}877 words}, cached as \texttt{data/cache/augmented\_word\_bank.json}.
The KNN retriever is otherwise unchanged; it just searches over a bigger universe.

\subsection*{Fix 2: 4-word-group hard-fail guard (extended \texttt{Deduplicator}).}

The rubric says: \emph{``If you generate a past connections puzzle you will automatically
fail.''} The original dedup only flagged 16-word-set overlap $> 6$ against past NYT
puzzles. We added:

\begin{itemize}[leftmargin=*, itemsep=0.1em]
\item A precomputed \texttt{\_nyt\_group\_set} of $\sim$2{,}216 past NYT 4-word groups
as frozensets, indexed at first call.
\item A new \texttt{check\_groups(groups)} method that flags if any generated 4-word
group matches a past NYT group exactly.
\item Inside CFR's retrieval, we now enumerate all $\binom{8}{4}=70$ combinations of
candidate words, rank them by avg pairwise cosine similarity, and return the first
combination that is \textbf{not} a verbatim NYT group.
\end{itemize}

Across 200 v2 puzzles (100 Mode A + 100 Mode B), \textbf{zero} accidentally reproduced a
past NYT group, confirming the guard works.

\subsection*{Verified impact}

\begin{table}[H]
\centering
\small
\begin{tabular}{l r r}
\toprule
\textbf{Metric} & \textbf{v1} & \textbf{v2} \\
\midrule
Word bank size                          & 4{,}918            & \textbf{14{,}877} \\
4-word group dedup                      & \ding{55}          & \textbf{\ding{51}} \\
Non-NYT words in output (Mode A)        & 0\%                & \textbf{62.6\%} \\
Non-NYT words in output (Mode B)        & 0\%                & \textbf{69.1\%} \\
Exact NYT group matches (200 puzzles)   & unknown            & \textbf{0} \\
Pass rate Mode A                        & 97\%               & \textbf{99\%} \\
Pass rate Mode B                        & 99\%               & \textbf{99\%} \\
\bottomrule
\end{tabular}
\end{table}

%-----------------------------------------------------------------------------
\section{Solver Benchmark (Sanity Check on NYT Data)}

To verify the validator is not over-accepting our own puzzles, we ran both solvers
on the 554 ground-truth NYT puzzles:

\begin{table}[H]
\centering
\small
\begin{tabular}{lrr}
\toprule
\textbf{Solver} & \textbf{Full-puzzle solve rate} & \textbf{Expected (paper baseline)} \\
\midrule
Embedding solver (greedy) & 2.0\% & \textasciitilde11.6\% (``Missed Connections'') \\
Clustering solver (beam)  & 3.4\% & Higher than embedding \\
\bottomrule
\end{tabular}
\caption{Solver performance on 554 real NYT puzzles.}
\end{table}

Our solvers correctly find only 2--3\% of real NYT puzzles (which often contain pop
culture / wordplay categories that defeat semantic similarity). This is the \emph{right}
kind of conservative validator --- a puzzle that \emph{our} solvers can crack
deterministically is one where the semantic structure is cleanly recoverable, i.e.,
a solvable puzzle.

%-----------------------------------------------------------------------------
\section{Deployed Artifacts}

\begin{itemize}[leftmargin=*, itemsep=0.15em]
\item \textbf{542-puzzle static web app}: live on GitHub Pages at
\url{https://kevin2330.github.io/nyt-connections-generator/}.
Vanilla JS frontend; client-side puzzle selection; no backend required.
\item \textbf{Color-split dataset}: 534 validated groups in each of 4 difficulty
buckets (yellow / green / blue / purple), totaling 8{,}544 NYT-style word groups.
\item \textbf{Jupyter notebook} (\texttt{notebooks/infinite\_connections\_demo.ipynb})
reproducing all generation and validation steps in dry-run mode.
\item \textbf{CFR module} (\texttt{src/cfr/}): 3 new Python files ($\sim$300 lines
total), completely separate from the baseline pipeline. All existing code is untouched.
\end{itemize}

%-----------------------------------------------------------------------------
\section{Summary of Current Results}

\begin{table}[H]
\centering
\small
\begin{tabular}{p{0.42\linewidth} p{0.52\linewidth}}
\toprule
\textbf{Artifact} & \textbf{Value} \\
\midrule
Total puzzles generated             & \textbf{1{,}175+} (775 baseline + 400 CFR) \\
Total validated puzzles             & \textbf{1{,}106+} (712 baseline + 394 CFR) \\
Best pass rate                      & \textbf{99\%} (CFR v2, both modes) \\
Fastest per-puzzle generation       & \textbf{1.27 s} (CFR v2 Mode A, 0 LLM calls) \\
Total spending to date              & $\sim$\$2.50 \\
Word bank (v2)                      & \textbf{14{,}877 words} (4{,}918 NYT + $\sim$10{,}000 WordNet) \\
Non-NYT word usage (v2)             & 62--69\% of output words \\
Past NYT group collisions (v2)      & \textbf{0 across 200 puzzles} (rubric safety) \\
Web app                             & Deployed, 542 puzzles playable online \\
Pipelines built                     & \textbf{2 novel} (CFR Mode A + Mode B); baseline kept as reference \\
Methodological novelty              & Embedding-first retrieval + group-level dedup \\
\bottomrule
\end{tabular}
\end{table}

\subsection*{Key Takeaway}

The project has produced both a competent LLM-heavy baseline \textbf{and} a new,
less-LLM-reliant methodology (CFR) that beats the baseline on every measurable axis:
pass rate, cost, and speed. CFR makes the academic contribution of the project
meaningfully original: we don't just wrap an LLM, we \textbf{use embeddings as
the primary generative substrate} and let the LLM play a small, well-defined role
(writing 4 category names). The system is fully reproducible, deployed, and
documented.

\end{document}
